\documentclass[a4paper]{article}
\usepackage{graphicx}
\usepackage{amsmath,amssymb}
\usepackage{hyperref}
\usepackage{minted}
\usepackage{multirow}

\title{anomaly\_vs\_outlier\_vs\_novelty}
\date{}

\begin{document}
    \maketitle
    \begin{itemize}
        \item \textbf{Outlier:}
        A data point that lies far from the majority of observations in feature space.
        It is a purely statistical notion (distance, density, deviation).
        An outlier is not necessarily an error; it may be rare but valid. (Such as noise)
        
        \item \textbf{Anomaly:}
        A behavior or state that violates the expected operating pattern of a system.
        It has semantic or system-level meaning (fault, attack, malfunction).
        An anomaly may appear as an outlier, but it can also be a structured pattern (e.g., temporal drift) rather than a single extreme point.
        
        \item \textbf{Novelty:}
        A sample or regime coming from a distribution not seen during training.
        Training data are assumed to contain only normal behavior.
        Novelty is typically interpreted as distribution shift rather than local abnormality.
    \end{itemize}
    
    % \bibliographystyle{plain}
    % \bibliography{/home/donkarlo/Dropbox/repo/shared_research/refs/refs}
\end{document}